\section{Our Approach}
\label{sec:impl}

All five implementations grow the same exact CART or C4.5 tree. They differ
in how the sorted order of a node's samples is obtained, in how many candidate
cuts are scored exactly, and in where and in what order the nodes are grown.
\ImplGpu{} and \ImplHybrid{} are two configurations of one CUDA code path.

\subsection{Common Design}
\label{sec:impl-common}

The training set is stored column-major as 32-bit floats. Every
implementation scores a cut with the same double-precision routine, compiled both for the CPU and for the GPU, and uses
the same rules to choose the attribute and to stop growth. Two gains within
$10^{-12}$ are treated as equal; the earlier cut, then the lower attribute
index, wins. The optimizations below only skip cuts that cannot win under this
rule, and the selected cut is always scored by the shared routine. The tree is
stored as flat arrays in preorder, so its layout does not depend on which
thread or device created a node. Our tests require byte-identical trees from
all implementations. Pruning runs afterwards on the CPU with shared code; the
experiments in this paper do not prune.

\subsection{\ImplSerialBase}
\label{sec:impl-simple}

\ImplSerialBase{} follows the textbook formulation and serves as the reference
for the others. At each node, it sorts the node's sample indices by every
attribute with a comparison sort, scores every admissible cut, and splits the
indices into two new lists, recursing on each. A
node with $n$ samples and $F$ attributes costs $O(Fn\log n)$, and the sort
reads values at random positions through the indices, so it is bound by
memory latency rather than bandwidth.

\subsection{\ImplSerialOpt}
\label{sec:impl-serial}

\ImplSerialOpt{} changes how the algorithm is executed in three ways.

\emph{Sort once.} As in SLIQ and SPRINT~\cite{mehta1996sliq,shafer1996sprint},
every attribute is sorted once before induction, here by a stable,
single-threaded radix sort on the bits of the values. A column consists of
8-byte entries that pair a value with its sample index and class, so a sweep
reads class counts without indirect accesses. A node occupies the same index
range in every column. After a split, the left samples of the winning
attribute are the first $n_L$ entries of its range; the other columns are
partitioned stably according to these samples, which keeps both children
sorted at a cost of $O(Fn)$ per node. An attribute that is constant over a
node is dropped for its whole subtree.

\emph{Score few cuts exactly.} The sweep updates a cheap estimate of the gain
at each cut, without divisions or a loop over the classes and with a bounded
rounding error. Only cuts whose estimate could still beat the best one are
scored by the shared routine, so the result equals that of scoring every cut.
For two classes, the gain is convex in the left class counts, and blocks of
cuts are skipped when a bound at the corners of their count range falls below
the best gain.

\emph{Allocate once.} The columns occupy $F+1$ slots of $8N$ bytes for $N$
training samples. The spare slot lets nodes of $2^{15}$ or more samples be
partitioned out of place with streaming stores; smaller nodes are partitioned
in place through a reused buffer. Large blocks are placed on transparent huge
pages. \ImplSerialOpt{} is the builder of \ImplParallel{} run without a
thread pool.

\subsection{\ImplParallel}
\label{sec:impl-parallel}

\ImplParallel{} runs this builder on a thread pool and, like the NAP strategy
of Aldinucci et al.~\cite{aldinucci2014fastflow}, chooses the kind of
parallelism by node size. Nodes of at least 32,768 samples sweep their
attributes in parallel and partition one column per thread. After a node of at
least 4,096 samples is split, its left child becomes a task and the thread
continues with the right child. Keeping $\min(F-1, p)$ spare slots for $p$
threads lets a large node move all its columns at once. The presort is
parallel too. Nodes are created concurrently and renumbered in preorder at the
end, so the tree equals that of \ImplSerialOpt{}.

The implementation has no explicit NUMA policy: it neither pins threads nor
binds or interleaves memory. The thread count follows the process's CPU
affinity, so \texttt{numactl} or \texttt{taskset} can confine a run to
selected sockets, and by default Linux places each page on the socket of the
thread that first writes it.

\subsection{\ImplGpu}
\label{sec:impl-gpu}

A small node cannot occupy thousands of GPU threads. \ImplGpu{} therefore grows
the tree breadth-first: all nodes of a level are processed by one fixed
sequence of kernel launches. Work is divided into tiles
of 2,048 entries of one node's range in one column, one thread block per tile.
The device holds one copy of the columns in $F+1$ slots, presorted with CUB's
stable radix sort as soon as each attribute is uploaded.

Within a warp, the class counts left of each entry are obtained with ballot
and population-count instructions rather than a serial scan. The best cut of
every node and attribute is copied to the host, which applies the same CART or
C4.5 rule as the CPU; this is the only round trip per level. A single-pass
stable partition with decoupled look-back~\cite{merrill2016scan} then writes
the children into a free slot and counts the class histograms of the next
level's tiles. The A100 and H100 score all cuts in double precision; GPUs with
slow double precision first discard hopeless cuts with a conservative
single-precision pass.

Data are uploaded through pinned staging buffers that CPU threads fill while
earlier chunks are in flight. This is the only explicit NUMA handling in the
code: on multi-socket hosts, the buffers are allocated on the GPU's NUMA
node and the copying threads are pinned to its CPUs. In \ImplGpu{}, the
hand-off threshold of \ImplHybrid{} is one sample, so every node stays on the
GPU; the CPU stages copies, chooses splits and, for C4.5, prunes.

\subsection{\ImplHybrid}
\label{sec:impl-hybrid}

\ImplHybrid{} divides the nodes of one tree between the devices by size
(Fig.~\ref{fig:hybrid}). A GPU level costs about as much for many small nodes as
for a few large ones, whereas CPU threads grow small subtrees independently.
Children with fewer than $m$ samples therefore leave the GPU and are grown by
the \ImplParallel{} builder while the GPU continues with the next level. By
default, $m$ is $N/8192$ bounded to $[1024, 2048]$ for CART and $N/4096$
bounded to $[1024, 4096]$ for C4.5, bounds chosen empirically.

The hand-off is asynchronous. A kernel gathers the handed-off nodes' columns,
except those of attributes constant on them, into a device ring buffer. Copy
streams move them through two pinned buffers into a pageable host arena, and a
host callback queued behind each copy submits the CPU tasks. The GPU never
waits for these tasks, and waits for the copies only when it is a whole ring
ahead of the PCIe link. The arena is pre-faulted by all threads in the
background and is not bound to the GPU's NUMA node. Nodes from both devices
enter one shared store, so the final tree is again the same.

\begin{figure}[t]
\centering
\includegraphics[width=\textwidth]{figs/approach-hybrid.pdf}
\caption{One level of \ImplHybrid{}. The GPU sweeps and partitions all large
nodes of the level together; the host chooses each node's split. Children
with fewer than $m$ samples are gathered into a device ring buffer, copied
through pinned buffers into a host arena, and grown as subtree tasks by the
CPU threads while the GPU continues with the next level. \ImplGpu{} is the
same scheme without the hand-off ($m=1$).}
\label{fig:hybrid}
\end{figure}
